% ---
% Dedicatória (opcional)
% ---
% \begin{dedicatoria}
%    \vspace*{\fill}
%    \centering
%    \noindent
%    \textit{[DRAFT --- dedicatória]} \vspace*{\fill}
% \end{dedicatoria}
% ---

% ---
% Agradecimentos (DESATIVADO nesta entrega — preencher antes da defesa)
% ---
% \begin{agradecimentos}
% [DRAFT --- agradecimentos a completar]
% \end{agradecimentos}
% ---

% ---
% Epígrafe
% ---
% \begin{epigrafe}
%     \vspace*{\fill}
% 	\begin{flushright}
% 		\textit{``[DRAFT --- epígrafe]''\\
% 		([DRAFT --- referência])}
% 	\end{flushright}
% \end{epigrafe}
% ---

% ---
% RESUMOS
% ---

% resumo em português
\setlength{\absparsep}{18pt}
\begin{resumo}
\SingleSpacing
O acervo cartográfico da Diretoria de Serviço Geográfico (DSG) é organizado por nomenclaturas técnicas --- Índice de Nomenclatura (INOM), Mapa Índice (MI), escalas e tipos de produto do Sistema Cartográfico Nacional --- cujo domínio é uma barreira para usuários não especialistas. O 1º Centro de Geoinformação (1º CGEO) desenvolveu um protótipo de busca em linguagem natural com Saída Estruturada e um único modelo local (Phi-4 14B), que demonstrou a viabilidade da abordagem, mas dependia de um dicionário de normalização fixo no código e não dispunha de um protocolo de avaliação. Este Projeto Final de Curso aprimorou a \textbf{camada de tradução linguística} dessa solução --- a etapa em que a consulta livre do usuário é convertida em parâmetros estruturados de busca. O \textit{backend} foi reimplementado em Python, com FastAPI e LangChain, e a Saída Estruturada foi substituída por \textit{Tool Calling} nativo, sem dicionário de normalização e sem exemplos no \textit{prompt}. Para medir a qualidade da tradução, construiu-se um \textit{dataset} de 310 consultas em três camadas --- casos de teste do protótipo, consultas redigidas pelos autores e consultas geradas programaticamente --- com gabarito definido por um manual de anotação, leituras múltiplas explícitas para as consultas ambíguas e uma auditoria em três etapas (checagens automáticas, anotação independente às cegas e adjudicação registrada), que obteve concordância quase perfeita entre o gabarito e a anotação independente. Um \textit{pipeline} reprodutível executou o \textit{dataset} sobre quatro modelos abertos executáveis localmente via Ollama (Qwen 3 4B, Gemma 4 E4B, Gemma 4 E2B e Mistral Nemo 12B) e, como referência, sobre \res{groq}{geral}{nmodelos} modelos de maior porte servidos em nuvem, medindo acurácia com intervalo de confiança, precisão, \textit{recall} e F1 por campo, latência e significância das diferenças entre modelos. Três dos modelos locais --- Gemma 4 E4B, Qwen 3 4B e Gemma 4 E2B --- tiveram acurácias estatisticamente equivalentes, de \res{principal}{e4b}{acuracia}, \res{principal}{qwen}{acuracia} e \res{principal}{e2b}{acuracia}; o Mistral Nemo 12B ficou em \res{principal}{mistral}{acuracia}, por descrever a chamada da ferramenta em texto em vez de emiti-la. Todos recusaram as consultas fora do domínio, e datas relativas e ordenação foram as categorias mais difíceis. Na estação de referência, com 4~GB de VRAM, o Gemma 4 E2B teve latência mediana de \res{estacao}{e2b}{latmeds}~s, contra \res{estacao}{qwen}{latmeds}~s do Qwen 3 4B, e é o modelo recomendado para uso operacional; os modelos em nuvem, nas mesmas consultas, acertaram \res{groqcomum}{gptoss120}{acuracia} e \res{groqcomum}{gptoss20}{acuracia}, o que indica a margem de melhoria disponível. O \textit{dataset} auditado e o \textit{pipeline} permitem repetir a comparação, com o mesmo critério, sempre que surgirem novos modelos.

\textbf{Palavras-chave}: Modelos de Linguagem de Grande Porte. \textit{Tool Calling}. Catálogo de Dados Espaciais. Geoinformação. Ollama. LangChain. Avaliação de Modelos.

\end{resumo}

% resumo em inglês
\begin{resumo}[Abstract]
 \begin{otherlanguage*}{english}
\SingleSpacing
The cartographic collection of the Brazilian Army Geographic Service Directorate (DSG) is organized by technical nomenclatures --- the Nomenclature Index (INOM), the Index Map (MI), and the scales and product types of the National Cartographic System --- whose mastery is a barrier for non-specialist users. The 1st Geoinformation Center (1º CGEO) developed a natural-language search prototype based on Structured Output and a single local model (Phi-4 14B), which demonstrated the feasibility of the approach but relied on a normalization dictionary hard-coded in the source and had no evaluation protocol. This Final Course Project improved the \textbf{natural-language translation layer} of that solution --- the step in which the user's free-form query is converted into structured search parameters. The backend was re-implemented in Python, with FastAPI and LangChain, and Structured Output was replaced by native Tool Calling, with no normalization dictionary and no worked examples in the prompt. To measure translation quality, a 310-query dataset was built in three layers --- the prototype's test cases, queries written by the authors, and programmatically generated queries --- with ground truth defined by a written annotation manual, explicit multiple readings for ambiguous queries, and a three-stage audit (automatic checks, blind independent annotation, and recorded adjudication) that showed near-perfect agreement between the ground truth and the independent annotation. A reproducible pipeline ran the dataset on four open models that run locally through Ollama (Qwen 3 4B, Gemma 4 E4B, Gemma 4 E2B, and Mistral Nemo 12B) and, as a reference, on \res{groq}{geral}{nmodelosen} larger cloud-served models, measuring accuracy with confidence intervals, per-field precision, recall and F1, latency, and the significance of differences between models. Three of the local models --- Gemma 4 E4B, Qwen 3 4B, and Gemma 4 E2B --- reached statistically equivalent accuracy, while Mistral Nemo 12B lagged far behind because it described the tool call in plain text instead of emitting it. All models correctly declined out-of-domain queries, and relative dates and sorting were the hardest categories. On the 4~GB reference workstation, Gemma 4 E2B had the lowest latency at the same accuracy and is the model recommended for operational use; the larger cloud-served models scored substantially higher on the same queries, which indicates the room for improvement available to local models. The audited dataset and the pipeline allow the comparison to be repeated, under the same criteria, whenever new models become available.

   \vspace{\onelineskip}
   \noindent
   \textbf{Keywords}: Large Language Models. Tool Calling. Spatial Data Catalog. Geoinformation. Ollama. LangChain. Model Evaluation.
 \end{otherlanguage*}
\end{resumo}
